\section*{Ethical Considerations}

\textbf{No real systems or users are attacked.} The main 30 scenarios
are synthetic: hand-authored specifications with LLM-generated surface
text (Section~\ref{sec:method}), not real production agent-memory
traces. The 20-scenario real-document validation slice
(Section~\ref{sec:real-document-slice}) uses source text copied
offline from public documents (NIST, CISA, OWASP, FTC, AWS, Azure,
FDA, IRS guidance and reference material); that text is modified only
within our own experimental corpus --- a fabricated claim is inserted
into our copy, never into the source document itself, and no live
system, deployed agent, real user account, or production document is
targeted, accessed, or compromised at any point in either corpus.

\textbf{Our poison-injection techniques are not novel attack
contributions.} The four \texttt{poison\_form} styles we use to
construct scenarios (\texttt{direct\_instruction}, \texttt{embedded\_fact},
\texttt{authoritative\_framing}, \texttt{multi\_hop\_setup}) are simple
content-injection patterns consistent with attack classes already
published and analyzed in prior work we cite (MPBench's attack
taxonomy, AgentPoison, MINJA). We do not contribute a new poisoning
technique, a new trigger-optimization method, or any capability that
lowers the bar for a real attacker beyond what is already public. The
contribution is a measurement framework and defensive metric, applied
against a known-compromise precondition (Section~\ref{sec:threat-model}).

\textbf{The primary contribution is defensive.} Blast-radius
reconstruction, attribution thresholding, and depth-aware containment
are all analysis and mitigation tools for an operator who has already
detected a compromise, not offensive techniques. To the extent our
results identify a weakness (e.g.\ the small but real laundering rate
in H3, or \texttt{CO\_RETRIEVED} under-detection in H1), they identify
weaknesses in \emph{defensive} provenance mechanisms, motivating
better defenses, not new attacks on real systems.

\textbf{Released artifacts do not increase attack capability.} The
scenario specifications, generation harness, and labeling pipeline we
release (Open Science) operate only on our own synthetic corpus and
require an attacker to already have the capability MemSecBench,
MemLineage, and MINJA already document as their starting threat model
(influence over content entering an agent's write path). Releasing our
harness does not provide a novel means of gaining that access.

\textbf{Dual-use consideration.} Any published work characterizing
where provenance-based defenses fail (H1's inflation growth, H2's
recall cost under thresholding, H3's laundering rate) could in
principle inform an adversary about which derivation transforms or
fan-out regimes are least likely to be caught. We judge the
operational disclosure risk here to be low relative to the defensive
value: our findings do not reveal a new attack, only quantify
existing, previously-known failure modes with more precision than
prior work. We are not aware of a responsible-disclosure obligation
here, since no deployed system or vendor-specific vulnerability is
implicated --- our findings are about the general mechanics of
derivation provenance, not a flaw in any named product.
